% Part: first-order-logic
% Chapter: syntax-and-semantics
% Section: terms-formulas

\documentclass[../../../include/open-logic-section]{subfiles}

\begin{document}

% SEGMENT OLP-0152-S01
\olfileid{fol}{syn}{frm}

\olsection{சொற்களும் \printtoken{P}{formula}s உம்}

ஒரு முதல்தர மொழி~$\Lang L$ கொடுக்கப்பட்டவுடன், $\Lang L$ இன் அடிப்படைச்
சொற்களஞ்சியத்திலிருந்து உருவாகும் வெளிப்பாடுகளை வரையறுக்கலாம். அவற்றில்
குறிப்பாக \emph{சொற்களும்} \emph{!!{formula}s} உம் அடங்கும்.

\begin{defn}[Terms]
\ollabel{defn:terms}
$\Lang L$ இன் \emph{சொற்களின்} தொகுப்பு~$\Trm[L]$ பின்வருமாறு
தொகுத்தறிதல் வழியில் வரையறுக்கப்படுகிறது:
\begin{enumerate}
\item ஒவ்வொரு !!{variable} உம் ஒரு சொல்.
\item $\Lang L$ இன் ஒவ்வொரு !!{constant} உம் ஒரு சொல்.
\item $f$ ஒரு $n$-இட !!{function} ஆகவும், $t_1$, \dots, $t_n$ சொற்களாகவும்
  இருந்தால், $\Atom{f}{t_1, \ldots, t_n}$ ஒரு சொல்.
\tagitem{limitClause}{வேறு எதுவும் சொல் அல்ல.}{}
\end{enumerate}
!!{variable}s எதுவும் இல்லாத சொல் \emph{மூடிய சொல்} எனப்படுகிறது.
\end{defn}

\begin{explain}
நமது மொழி வரையறையிலும் சொல் வரையறையிலும் !!{constant}s தனிப் பிரிவு
குறிகளாகத் தோன்றுகின்றன; ஆனால் அவற்றை பூஜ்ய-இட !!{function}s ஆகவும்
சேர்க்கலாம். அப்படியானால் சொல் வரையறையின் இரண்டாவது விதி தேவையில்லை.
$n=0$ ஆக இருக்கும்போது $\Atom{f}{t_1, \ldots, t_n}$ ஐத் தனித்த
$f$ ஆகப் புரிந்துகொண்டால் போதும்.
\end{explain}

\begin{defn}[Formulas]
\ollabel{defn:formulas}
மொழி~$\Lang L$ இன் \emph{!!{formula}s} தொகுப்பு~$\Frm[L]$ பின்வருமாறு
தொகுத்தறிதல் வழியில் வரையறுக்கப்படுகிறது:
\begin{enumerate}
\tagitem{prvFalse}{$\lfalse$ ஒரு அணு !!{formula}.}{}

\tagitem{prvTrue}{$\ltrue$ ஒரு அணு !!{formula}.}{}

\item $R$ என்பது $\Lang L$ இன் $n$-இட !!{predicate} ஆகவும், $t_1$, \dots,
  $t_n$ என்பது $\Lang L$ இன் சொற்களாகவும் இருந்தால்,
  $\Atom{R}{t_1,\ldots, t_n}$ ஒரு அணு !!{formula}.

\item $t_1$, $t_2$ ஆகியவை $\Lang L$ இன் சொற்களாக இருந்தால்,
  $\Atom{\eq}{t_1, t_2}$ ஒரு அணு !!{formula}.

\tagitem{prvNot}{$!A$ ஒரு !!a{formula} என்றால், $\lnot !A$ ஒரு
  !!a{formula}.}{}

\tagitem{prvAnd}{$!A$, $!B$ ஆகியவை !!{formula}s என்றால்,
  $(!A \land !B)$ ஒரு !!a{formula}.}{}

\tagitem{prvOr}{$!A$, $!B$ ஆகியவை !!{formula}s என்றால்,
  $(!A \lor !B)$ ஒரு !!a{formula}.}{}

\tagitem{prvIf}{$!A$, $!B$ ஆகியவை !!{formula}s என்றால்,
  $(!A \lif !B)$ ஒரு !!a{formula}.}{}

\tagitem{prvIff}{$!A$, $!B$ ஆகியவை !!{formula}s என்றால்,
  $(!A \liff !B)$ ஒரு !!a{formula}.}{}

\tagitem{prvAll}{$!A$ ஒரு !!a{formula} ஆகவும் $x$ ஒரு !!a{variable} ஆகவும்
  இருந்தால், $\lforall[x][!A]$ ஒரு !!a{formula}.}{}

\tagitem{prvEx}{$!A$ ஒரு !!a{formula} ஆகவும் $x$ ஒரு !!a{variable} ஆகவும்
  இருந்தால், $\lexists[x][!A]$ ஒரு !!a{formula}.}{}

\tagitem{limitClause}{வேறு எதுவும் !!a{formula} அல்ல.}{}
\end{enumerate}
\end{defn}

\begin{explain}
சொற்களின் தொகுப்புக்கும் !!{formula}s தொகுப்புக்கும் அளித்த வரையறைகள்
\emph{தொகுத்தறிதல் வரையறைகள்}. அடிப்படையில், !!{formula}s தொகுப்பை எண்ணற்ற
நிலைகளில் உருவாக்குகிறோம். ஆரம்ப நிலையில் எல்லா அணு சூத்திரங்களையும்
சூத்திரங்களாக அறிவிக்கிறோம்; இது வரையறையின் முதல் சில நிகழ்வுகளான
\iftag{prvTrue}{$\ltrue$, }{}%
\iftag{prvFalse}{$\lfalse$, }{}%
$\Atom{R}{t_1,\dots,t_n}$ மற்றும் $\Atom{\eq}{t_1,t_2}$ ஆகியவற்றுக்கு
ஒத்ததாகும். எனவே ``அணு !!{formula}'' என்பது இந்த வடிவிலுள்ள எந்த !!{formula}
என்றும் பொருள்.

வரையறையின் பிற நிகழ்வுகள் ஏற்கெனவே உருவாக்கப்பட்ட !!{formula}s இலிருந்து
புதிய !!{formula}s ஐ உருவாக்கும் விதிகளைத் தருகின்றன. இரண்டாவது நிலையில்
அணு சூத்திரங்களிலிருந்து புதிய !!{formula}s ஐ உருவாக்கலாம். மூன்றாவது
நிலையில் அணு !!{formula}s மற்றும் இரண்டாவது நிலையில் உருவானவற்றிலிருந்து
புதிய சூத்திரங்களை உருவாக்கலாம்; இவ்வாறு தொடர்கிறது. அத்தகைய ஒரு நிலையில்
இறுதியில் உருவாக்கப்படும், வேறு எதுவும் அல்லாத ஒவ்வொன்றும் ஒரு !!{formula}.
\end{explain}

மரபுப்படி, $\eq$ ஐ அதன் இரு உறுப்புகளுக்கு இடையில் எழுதி அடைப்புக்குறிகளை
விடுகிறோம்: $\eq[t_1][t_2]$ என்பது $\Atom{\eq}{t_1,t_2}$ க்கான சுருக்கம்.
மேலும், $\lnot \Atom{\eq}{t_1,t_2}$ என்பது $\eq/[t_1][t_2]$ எனச்
சுருக்கப்படுகிறது. $!B$, $!C$ ஆகியவற்றிலிருந்து இரு-இடத் தொடுப்புக்குறி
$\ast$ கொண்டு உருவாக்கப்படும் $(!B \ast !C)$ வடிவில் வெளிப்புற அடைப்புக்குறிகளை
பல நேரம் விட்டு, வெறுமனே~$!B \ast !C$ என்று எழுதுவோம்.

\begin{intro}
சில தருக்க நூல்கள், !!{variable}~$x$ ஆனது $!A$ இல் தோன்ற வேண்டும் என்று
வலியுறுத்துகின்றன; அப்போதுதான்
\iftag{prvEx}{$\lexists[x][!A]$ }{}%
\iftag{notprvEx,notprvAll}{}{மற்றும் }%
\iftag{prvAll}{$\lforall[x][!A]$ }{}%
ஆகியவை
\iftag{notprvEx,notprvAll}{!!a{formula}}{!!{formula}s} எனக் கருதப்படும்.
இந்த நிபந்தனையை விதிக்காவிட்டாலும் எந்தத் தீங்கும் இல்லை; அது வேலைகளை
எளிதாக்குகிறது.
\end{intro}

\begin{tagblock}{defNot,defOr,defAnd,defIf,defIff,defTrue,defFalse,defEx,defAll}
\begin{defn}
வரையறுக்கப்பட்ட இயக்கிகளைக் கொண்டு உருவாக்கப்படும் சூத்திரங்கள் பின்வருமாறு
புரிந்துகொள்ளப்பட வேண்டும்:

\begin{tagenumerate}{defTrue,defFalse,defNot,defOr,defAnd,defIf,defIff,defEx,defAll}

\tagitem{defTrue}{$\ltrue$ என்பது
  \iftag{prvFalse}{$\lnot\lfalse$}{$(!A \lor \lnot !A)$ இன் சுருக்கம்; இங்கு
    $!A$ ஒரு நிலையான அணு !!{formula}.}}{}

\tagitem{defFalse}{$\lfalse$ என்பது
  \iftag{prvTrue}{$\lnot\ltrue$}{$(!A \land \lnot !A)$ இன் சுருக்கம்; இங்கு
    $!A$ ஒரு நிலையான அணு !!{formula}.}}{}

\tagitem{defNot}{$\lnot !A$ என்பது $!A \lif \lfalse$ என்பதன் சுருக்கம்.}{}

\tagitem{defOr}{$!A \lor !B$ என்பது
  \iftag{prvAnd}{$\lnot(\lnot !A \land \lnot !B)$}{$\lnot !A \lif
    !B$} என்பதன் சுருக்கம்.}{}

\tagitem{defAnd}{$!A \land !B$ என்பது
  \iftag{prvOr}{$\lnot(\lnot !A \lor \lnot !B)$}{$\lnot (!A \lif
    \lnot !B)$} என்பதன் சுருக்கம்.}{}

\tagitem{defIf}{$!A \lif !B$ என்பது
  \iftag{prvOr}{$\lnot !A \lor !B)$}{$\lnot (!A \land \lnot !B)$} என்பதன்
  சுருக்கம்.}{}

\tagitem{defIff}{$!A \liff !B$ என்பது $(!A \lif !B) \land (!B
  \lif !A)$ என்பதன் சுருக்கம்.}{}

\tagitem{defAll}{$\lforall[x][!A]$ என்பது
  $\lnot\lexists[x][\lnot !A]$ என்பதன் சுருக்கம்.}{}

\tagitem{defEx}{$\lexists[x][!A]$ என்பது
  $\lnot\lforall[x][\lnot !A]$ என்பதன் சுருக்கம்.}{}
\end{tagenumerate}
\end{defn}
\end{tagblock}

குறிப்பிட்ட பயன்பாட்டுக்கான மொழியில் வேலை செய்தால், இரு-இட !!{predicate}s
மற்றும் !!{function}s ஐ அவற்றின் சொற்களுக்கு இடையில் எழுதுவோம்; உதாரணமாக,
கணித மொழியில் $t_1 < t_2$, $(t_1 + t_2)$ என்றும், கணக்குத்தொகுப்புக்
கோட்பாட்டின் மொழியில் $t_1 \in t_2$ என்றும் எழுதலாம். கணித மொழியின்
வாரிசுச் சார்பு வழக்கமாக அதன் வாதத்திற்குப் \emph{பிறகு}:~$t'$ என்று
எழுதப்படுகிறது. அதிகாரப்பூர்வமாக, இவை முறையே
$\Atom{\Obj A^2_0}{t_1, t_2}$, $\Obj f^2_0(t_1, t_2)$,
$\Atom{\Obj A^2_0}{t_1, t_2}$ மற்றும் $f^1_0(t)$ ஆகியவற்றுக்கான
மரபுச் சுருக்கங்கள்.

\begin{defn}[Syntactic identity]
$\ident$ குறி குறியீடுகளின் சரங்களுக்கிடையிலான தொடரமைப்பு அடையாளத்தை
வெளிப்படுத்துகிறது; அதாவது, $!A \ident !B$ என்பது $!A$, $!B$ இரண்டும் ஒரே
நீளமுள்ள குறியீட்டு சரங்களாக இருந்து ஒவ்வொரு இடத்திலும் ஒரே குறியைக்
கொண்டிருக்கும்போது, அப்போதும் அப்போது மட்டுமே.
\end{defn}

$\ident$ குறியின் இருபுறமும் இணைப்பால் பெறப்பட்ட சரங்கள் இருக்கலாம்.
உதாரணமாக, $!A \ident (!B \lor !C)$ என்றால்: $!A$ என்ற குறியீட்டு சரம்,
திறப்பு அடைப்புக்குறி, $!B$ சரம், $\lor$ குறி, $!C$ சரம், மூடு
அடைப்புக்குறி ஆகியவற்றை இந்த வரிசையில் இணைத்துப் பெறப்படும் சரத்துடன்
ஒன்றே. இந்நிலையில் $!A$ இன் முதல் குறி திறப்பு அடைப்புக்குறி என்றும்,
இரண்டாவது குறியிலிருந்து தொடங்கி $!A$ இல் $!B$ துணைச்சரம் இருப்பதாகவும்,
அதனைத் தொடர்ந்து $\lor$ வருவதாகவும் நமக்குத் தெரியும்.

சொற்களும் !!{formula}s உம் தொகுத்தறிதல் வரையறைகளால் அடிப்படை உறுப்புகளிலிருந்து
உருவாக்கப்படுவதால், அவற்றைப் பற்றிய கூற்றுகளை நிறுவப் பின்வரும்
தொகுத்தறிதல் கொள்கைகளைப் பயன்படுத்தலாம்.

\begin{lem}[\emph{Principle of induction on terms}]
    \ollabel{lem:trmind}
    $\Lang L$ ஒரு முதல்தர மொழியாக இருக்கட்டும்.
    ஒரு பண்பு~$P$ பின்வருமாறு இருந்தால்:
    %
    \begin{enumerate}
        \item அது ஒவ்வொரு !!{variable}~$v$ க்கும் பொருந்தும்,
        %
        \item அது $\Lang L$ இன் ஒவ்வொரு !!{constant}~$a$ க்கும் பொருந்தும்,
        மற்றும்
        %
        \item $f$ என்பது $\Lang L$ இன் $n$-இட !!{function} ஆகவும்,
            $t_1$,~\dots, $t_n$ க்குப் பண்பு பொருந்தும்போதும்
            $f(t_1,\dotsc,t_n)$ க்குப் பொருந்தும்
    \end{enumerate}
    ($t_1$,~\dots, $t_n$ ஆகியவை $\Lang{L}$ இன் சொற்கள் என்ற நிபந்தனையில்),
    அப்போது $\Trm[L]$ இன் ஒவ்வொரு சொல்லுக்கும்~$P$ பொருந்தும்.
\end{lem}

\begin{prob}
    \olref[fol][syn][frm]{lem:trmind} ஐ நிறுவுக.
\end{prob}

\begin{lem}[\emph{Principle of induction on !!{formula}s}]
    \ollabel{thm:frmind}
    $\Lang L$ ஒரு முதல்தர மொழியாக இருக்கட்டும்.
    சில பண்பு~$P$ எல்லா அணு !!{formula}s க்கும் பொருந்தி,
    %
    \begin{enumerate}
        \tagitem{prvNot}{$!A$ க்குப் பொருந்தும்போது $\lnot !A$ க்கும்
            பொருந்தும்;}{}
        \tagitem{prvAnd}{$!A$, $!B$ இரண்டுக்கும் பொருந்தும்போது
            $(!A \land !B)$ க்கும் பொருந்தும்;}{}
        \tagitem{prvOr}{$!A$, $!B$ இரண்டுக்கும் பொருந்தும்போது
            $(!A \lor !B)$ க்கும் பொருந்தும்;}{}
        \tagitem{prvIf}{$!A$, $!B$ இரண்டுக்கும் பொருந்தும்போது
            $(!A \lif !B)$ க்கும் பொருந்தும்;}{}
        \tagitem{prvIff}{$!A$, $!B$ இரண்டுக்கும் பொருந்தும்போது
            $(!A \liff !B)$ க்கும் பொருந்தும்;}{}
        \tagitem{prvEx}{$!A$ க்குப் பொருந்தும்போது
            $\lexists[x][!A]$ க்கும் பொருந்தும்;}{}
        \tagitem{prvAll}{$!A$ க்குப் பொருந்தும்போது
            $\lforall[x][!A]$ க்கும் பொருந்தும்;}{}
  \end{enumerate}
  ($!A$, $!B$ ஆகியவை $\Lang{L}$ இன் !!{formula}s என்ற நிபந்தனையில்),
  அப்போது $\Frm[L]$ இன் எல்லா சூத்திரங்களுக்கும்~$P$ பொருந்தும்.
\end{lem}

\end{document}
